El siguiente código de MATLAB representa un generador de números aleatorios donde ninguno de ellos se repite:

\begin{lstlisting}
clear all
close all
clc

a = randi([1, 15], 1, 15);
S = [];

while true
    for i = a
        if ismember(i, S)
            n = length(S);
            a = randi([1, 15], 1, 15 - n);
            break;
        end
        S = [S, i];
    end
    m = length(S);
    if m == 15
        break;
    end
end

S 
>> 15 12 10 8 4 5 9 7 3 13 14 1 2 11 6
\end{lstlisting}


Básicamente lo que se hizo es crear un arreglo de números aleatorios, es decir, la variable \texttt{a}. Fue necesario crear otra variable, \texttt{S}, donde guardaremos todos los números aleatorios que no se repiten. Al ponerlo dentro de un bucle \textsl{while} éste proceso se repetirá hasta que \texttt{S} no tenga ningun elemento repetido. \texttt{ismember} es una función que nos dice si dicho elemento está en \texttt{S} o no. En el caso que esté, se creará un nuevo arreglo pero en lugar de 15 será un arreglo de 15 - \texttt{n}, donde \texttt{n} es el tamaño de \texttt{S}. Cabe aclarar que este código se hizo sin ayuda de IA. 

Existe una función interna de MATLAB que hace lo mismo y es \texttt{randperm}:

\begin{lstlisting}
r = randperm(15)
>> 4 10 5 14 13 9 15 2 6 8 12 11 7 1 3
\end{lstlisting}


En el anterior código se escribió lo siguiente:

\begin{lstlisting}
S = [];
.
.
.
S = [S, i];
\end{lstlisting}


La variable \texttt{S} no está definida en tamaño, por lo tanto su tamaño irá cambiando a medida que le agreguemos elementos. MATLAB advierte lo siguiente:

\begin{quote}
    Variable appears to change size on very loop iteration. Consider preallocating for speed.
\end{quote}

Básicamente hay que definir desde un principio el tamaño del arreglo. Veamos los siguientes 2 ejemplos

\begin{lstlisting}
tic
v = [];
i = 1;
while i<sqrt(7777777)
 v = [ v i^2 ];
 i = i + 1;
end
toc
>> Elapsed time is 0.008991 seconds.


tic
n = floor(sqrt(7777777));
v = nan(1, n);
i = 1;
while i<sqrt(7777777)
 v(i) = i^2;
 i = i + 1;
end
toc
>> Elapsed time is 0.005196 seconds.

\end{lstlisting}


En el primer caso definimos un arreglo vacío y en el segundo predefinimos el tamaño que tendrá la lista. Se observa que el segundo ejemplo fue 57.79\% más rápido que el primer ejemplo. Por lo tanto, definir el tamaño del arreglo resulta una ventaja en el tiempo de compilación.






